\section{Some probabilistic claims}\label{app:prob}

%\commentt{check that we use all claims in the end}
% 9.7: checked that first claim and corollary are used (the second claim is used for the corollary

\begin{claim}\label{clm:var1}
For integers $s_1$ and $s_2$  let
$\{\chi_{i,j}\}_{(i,j)\in [s_1]\times [s_2]}$ be Bernoulli random variables for which the following holds.
\begin{enumerate}
\item  $\Pr[\chi_{i,j} = 1] = \mu$ for every $(i,j)\in [s_1]\times [s_2]$.
\item For every $i_1,i_2 \in [s_1]$, $j_1,j_2\in [s_2]$ such that $i_1\neq i_2$ and $j_1\neq j_2$,
the random variables $\chi_{i_1,j_2}$ and $\chi_{i_2,j_2}$ are independent.
\item For every $i\in [s_1]$, $j_1,j_2\in [s_2]$ such that  $j_1\neq j_2$,
$\Pr[\chi_{i,j_1}=\chi_{i,j_2} = 1] = \mu_{1,2}$, and for every $i_1,i_2 \in [s_1]$ and $j \in [s_2]$ such that $i_1\neq i_2$,  $\Pr[\chi_{i_1,j}=\chi_{i_2,j} = 1] = \mu_{2,1}$.
\end{enumerate}
Then
$$\Var\left[\sum_{i\in [s_1],j\in [s_2]}\chi_{i,j}\right] \leq s_1\cdot s_2 \cdot \mu + s_1\cdot {s_2\choose 2}\cdot \mu_{1,2} + {s_1\choose 2} \cdot s_2\cdot \mu_{2,1}\;.$$
\end{claim}

\begin{proof}
For each pair $(i,j)\in [s_1]\times [s_2]$,
let $\bar{\chi}_{i,j} = \chi_{i,j} - \mu$, so that $\bar{\chi}_{i,j} \leq \chi_{i,j}$, and
$\Ex[\bar{\chi}_{i,j}] = 0$. By the definition of the variance,
\begin{equation}\label{eq:var1}
\Var\left[\sum_{i\in [s_1],j\in [s_2]}\chi_{i,j}\right] =
  \sum_{(i_1,j_2),(i_2,j_2)\in [s_1]\times[s_2]} \Ex[\bar{\chi}_{i_1,j_1}\cdot \bar{\chi}_{i_2,j_2}]\;.
\end{equation}

We upper bound the right-hand-side of Equation~\eqref{eq:var1} by breaking
 the sum $\sum_{i_1,j_2,i_2,j_2} \Ex[\bar{\chi}_{i_1,j_1}\cdot \bar{\chi}_{i_2,j_2}]$ into four sub-sums.

 \paragraph{The first} is over $(i_1,j_1), (i_2,j_2)\in [s_1]\times [s_2]$ such that all four indices are distinct. By the second condition in the claim, for such quadruples of indices, $\chi_{i_1,j_1}$ is independent of $\chi_{i_2,j_2}$ and so
\begin{equation}\label{eq:cheb-case1}
\Ex[\bar{\chi}_{i_1,j_1}\cdot \bar{\chi}_{i_2,j_2}] =  \Ex[\bar{\chi}_{i_1,j_1}]\cdot \Ex[\bar{\chi}_{i_2,j_2}] = 0\;.
\end{equation}

\paragraph{The second} is over $i_1 = i_2$ and $j_1 = j_2$, so that
\begin{equation}\label{eq:cheb-case2}
\Ex[\bar{\chi}_{i_1,j_1}\cdot \bar{\chi}_{i_2,j_2}] = \Ex[(\bar{\chi}_{i_1,j_1})^2]
\leq \Ex[(\chi_{i_1,j_1})^2] = \Ex[\chi_{i_1,j_1}] = \mu \;,
\end{equation}
where the second inequality follows from the definition of $\bar{\chi}_{i,j}$.
There are $s_1\times s_2$ such quadruples.

\paragraph{The third} is over $i_1 = i_2$ and $j_1 \neq j_2$, in which case we have
\begin{equation}\label{eq:cheb-case3}
%\Ex[\bar{\chi}_{i_1,j_1}\cdot \bar{\chi}_{i_2,j_2}]  =
\Ex[\bar{\chi}_{i_1,j_1}\cdot \bar{\chi}_{i_1,j_2}]
\leq \Ex[\chi_{i_1,j_1}\cdot \chi_{i_1,j_2}] = \Pr[\chi_{i_1,j_1}=\chi_{i_1,j_2}=1]
  = \mu_{1,2}\;.
\end{equation}
There are $s_1\times {s_2\choose 2}$ such pairs.

\paragraph{The fourth} is over $i_1 \neq i_2$ and $j_1 = j_2$, where similarly to the third case we get
\begin{equation}\label{eq:cheb-case4}
\Ex[\bar{\chi}_{i_1,j_1}\cdot \bar{\chi}_{i_2,j_2}]  %= \Ex[\bar{\chi}_{i_1,j_1}\cdot \bar{\chi}_{i_2,j_1}]
\leq %\Ex[\chi_{i_1,j_1}\cdot \chi_{i_2,j_1}] = \Pr[\chi_{i_1,j_1}=\chi_{i_1,j_2}]
\mu_{2,1}\;
\end{equation}
where there are ${s_1\choose 2}\cdot s_2$ such pairs.

The claim follows by summing all four sub-sums.
\end{proof}


The next claim is very similar to Claim~\ref{clm:var1}, but the setting is slightly different, and will be used in our analysis several time.
In what follows, for an integer $s$ we use $\Phi(s)$ to denote the set of pairs $i,j\in [s]$ such that $i<j$.
\begin{claim}\label{clm:var2}
For an integer $s$l let $\{\chi_{i,j}\}_{(i,j)\in \Phi(s)}$ be Bernoulli random variables for which the following holds.
\begin{enumerate}
\item  $\Pr[\chi_{i,j} = 1] = \mu$ for every $(i,j)\in \Phi(s)$.
\item For every $i_1,i_2 \in [s_1]$, $j_1,j_2\in [s_2]$ such that $i_1\neq i_2$ and $j_1\neq j_2$,
the random variables $\chi_{i_1,j_2}$ and $\chi_{i_2,j_2}$ are independent.
\item For every $(i_1,j_1)\in \Phi(s)$ and $(i_2,j_2)\in \Phi(s)$  such that exactly two of the four indices are the same,
% $|\{i_1,j_1,i_2,j_2\}|=3$,
$\Pr[\chi_{i_1,j_1}=\chi_{i_2,j_2} = 1] = \mu_{1,2}$.
\end{enumerate}
Then $\Var\left[\sum_{(i,j)\in \Phi(s)}\chi_{i,j}\right] \leq {s\choose 2}\cdot \mu + {s\choose 3}\mu_{1,2}$.
\end{claim}
The proof of Claim~\ref{clm:var2} is very similar to the proof of Claim~\ref{clm:var1}, and is hence omitted. As a corollary of this claim we get:

\begin{corollary}\label{cor:cheb1}
For an integer $s$
let $\{\chi_{i,j}\}_{(i,j)\in \Phi(s)}$ be Bernoulli random variables where
 $\Pr[\chi_{i,j} = 1] = \mu$ for every $(i,j)\in \Phi(s)$. Suppose that the following conditions hold
 for some $c_1>0$ and $c_2 > 4$.
\begin{enumerate}
\item For every $(i_1,j_1)\in \Phi(s)$ and $(i_2,j_2)\in \Phi(s)$ such that the four indices are distinct, $\chi_{i_1,j_1}$
% \talya{changed from $\chi_{i_1, j_2}$ to $\chi_{i_1, j_1}$} 
and $\chi_{i_2,j_2}$ are independent.
\item For every $(i_1,j_1)\in \Phi(s)$ and $(i_2,j_2)\in \Phi(s)$  such that exactly two of the four indices are the same,
% $|\{i_1,j_1,i_2,j_2\}|=3$,
$\Pr[\chi_{i_1,j_1}=\chi_{i_2,j_2} = 1] \leq c_1\cdot \mu^{3/2}$.
\item % For every $(i,j)\in \Phi(s)$, $\Pr[\chi_{i,j} = 1] = \mu$, where $\mu \geq \frac{1}{c s^2}$.
   $s \geq c_2/\sqrt{\mu}$\;.
\end{enumerate}
Then $\Pr\left[ \sum_{(i,j)\in \Phi(s)} \chi_{i,j} = 0\right] \leq \frac{1+c_1}{c_2}$.
\end{corollary}

\begin{proof}
By Chebishev's inequality,
\begin{equation}\label{eq:cheb1}
\Pr\left[\sum_{(i,j)\in \Phi(s)}\chi_{i,j}=0\right] \leq \frac{\Var\left[\sum_{(i,j)\in \Phi(s)} \chi_{i,j}\right]}{\left(\Ex\left[\sum_{(i,j)\in \Phi(s)}\chi_{i,j}\right]\right)^2}\;
% = \frac{\sum_{i_1,j_2,i_2,j_2} \Ex[\bar{\chi}_{i_1,j_1}\cdot \bar{\chi}_{i_2,j_2}]}{{s\choose 2}^2\cdot \mu^2} \;.
\end{equation}
% we detect a triangle when selecting $\theta$ random neighbors of $a$ and $\theta$ random neighbors of $b$.
The denominator of the right-hand-side of Equation~\eqref{eq:cheb1} equals
${s\choose 2}^2\cdot \mu^2$.
By Claim~\ref{clm:var2} and the first two conditions of the current claim, the numerator of the right-hand-side of Equation~\eqref{eq:cheb1} is upper bounded by
${s\choose 2}\cdot \mu + {s\choose 3}\cdot c_1\cdot \mu^{3/2}$, so that the claim follows.
\end{proof}